
\section{引言}
\subsection{问题背景}

生鲜商超中蔬菜类商品保鲜期短，品相随销售时间增加而变差，大部分品种当日未售出则隔日无法再售。商超需根据历史销售和需求情况每日补货，但进货交易时间在凌晨3:00--4:00，商家须在不确切知道具体单品和进货价格的情况下做出补货决策。蔬菜的定价一般采用成本加成定价方法，商超对运损和品相变差的商品通常进行打折销售。可靠的市场需求分析对补货决策和定价决策尤为重要。从需求侧看，蔬菜类商品销售量与时间存在关联关系；从供给侧看，蔬菜供应品种在4月至10月较为丰富，商超销售空间的限制使得合理的销售组合极为重要。

题目提供某商超六个蔬菜品类（花叶类、辣椒类、食用菌、水生根茎类、茄类、花菜类）共251个单品的数据：附件1为商品信息，附件2为2020年7月1日至2023年6月30日的销售流水明细（878,503条交易记录），附件3为同期批发价格数据，附件4为近期损耗率数据。本文仅使用附件4中的可见工作表（单品级损耗率），品类级损耗率通过附件1与附件4的关联自行计算。

\subsection{问题重述}

\textbf{问题1：} 分析蔬菜各品类及单品销售量的分布规律及相互关系。

\textbf{问题2：} 以品类为单位，分析销售总量与成本加成定价的关系，给出各品类2023年7月1--7日的日补货总量和定价策略，使商超收益最大。

\textbf{问题3：} 在销售空间有限的情况下制定单品补货计划，可售单品总数27--33个，各单品订购量$\geq 2.5$ kg。根据6月24--30日可售品种，给出7月1日的单品补货量和定价策略，在尽量满足市场对各品类需求的前提下使收益最大。

\textbf{问题4：} 建议商超还需采集哪些数据，说明这些数据对解决上述问题的帮助和理由。
